	%%----------------------------------------------------------- TITLE COVER   --------------------------------------------------------------
\thispagestyle{empty}
\begin{minipage}{0.15\linewidth}
	\begin{figure}[H]
		\includegraphics[width=3cm]{images/facsciences.jpg}
	\end{figure}
\end{minipage}
\centering
\begin{minipage}{0.6\linewidth}
	\centering
	\large\textsc{UNIVERSITÉ D’ ANTANANARIVO\\[-3mm]FACULTÉ DES SCIENCES}
\end{minipage}
\begin{minipage}{0.15\linewidth}
	\begin{figure}[H]
		\begin{flushright}
			\includegraphics[width=3cm]{images/univ.jpg}
		\end{flushright}
	\end{figure}
\end{minipage}
\vspace{-1cm}
\begin{figure}[H]
	\centering
	\includegraphics[width=3cm]{images/physique.png}
\end{figure}

\begin{table}[H]
	\centering
	\begin{tabular}{rl}
		\textbf{Domaine :}& SCIENCES ET TECHNOLOGIES\\
		\textbf{Mention :}& PHYSIQUE ET APPLICATIONS\\
		\textbf{Parcours :}& ASTROPHYSIQUE ET ASTRONOMIE\\
	\end{tabular}
\end{table}

\begin{center}
	\normalsize{M\'EMOIRE en vue de l'obtention du diplôme de :}\\[3pt]
	\Large\textsc{MASTER II EN PHYSIQUE ET APPLICATIONS}\\[4pt]
	\normalsize{Intitulé :}\\[18pt]
		\colorbox{blue!20}{%
\parbox{\linewidth}{\centering\fontsize{18}{22}\selectfont\linespread{1}\selectfont\bfseries La répartition de l'hydrogène atomique (\HI) cinématiquement froid et chaud et son lien avec la formation stellaire dans les galaxies proches}}
	\\[25pt]
	\normalsize{Présenté le 12 Novembre 2026 par:}\\
	\large\textbf{Anjia Randia RABARISOA}\\[2em]
	%\large\textbf{Cirno}\\[30pt]
\end{center}
\noindent \flushleft Devant la commission d'examen composée par :

\begin{center}
	\begin{tabular*}{\linewidth}{@{\extracolsep{\fill}}lll}
		Président :    & \textbf{M. Minoson RAKOTOMALALA}        & Professeur Titulaire\\
		Rapporteur :   & \textbf{M. Roger IANJAMASIMANANA}       & Chercheur Postdoctoral\\
		Examinateurs : & \textbf{M. Solohery RANDRIAMAMPANDRY}   & Professeur\\
		& \textbf{M. Toky RANDRIAMAMPANDRY}       & Maître de Conférences\\
	\end{tabular*}
\end{center}
%-------------------------------------------------------------------------------------============================================
\newpage
\thispagestyle{empty}
\begin{minipage}{0.15\linewidth}
	\begin{figure}[H]
		\includegraphics[width=3cm]{images/facsciences.jpg}
	\end{figure}
\end{minipage}
\centering
\begin{minipage}{0.6\linewidth}
	\centering
	\large\textsc{UNIVERSITÉ D’ ANTANANARIVO\\[-3mm]FACULTÉ DES SCIENCES}
\end{minipage}
\begin{minipage}{0.15\linewidth}
	\begin{figure}[H]
		\begin{flushright}
			\includegraphics[width=3cm]{images/univ.jpg}
		\end{flushright}
	\end{figure}
\end{minipage}
\vspace{-3em}
\begin{figure}[H]
	\centering
	\includegraphics[width=3cm]{images/physique.png}
\end{figure}
\vspace{-1em}

\begin{table}[H]
	\centering
	\begin{tabular}{rl}
		\textbf{Domaine :}& SCIENCES ET TECHNOLOGIES\\
		\textbf{Mention :}& PHYSIQUE ET APPLICATIONS\\
		\textbf{Parcours :}& ASTROPHYSIQUE ET ASTRONOMIE\\[-1em]
	\end{tabular}
\end{table}

\begin{center}
	\normalsize{M\'EMOIRE en vue de l'obtention du diplôme de :}\\
	\Large\textsc{MASTER II EN PHYSIQUE ET APPLICATIONS}\\[3pt]
	\normalsize{Intitulé :}\\[1em]
		\colorbox{blue!20}{%
\parbox{\linewidth}{\centering\fontsize{18}{22}\selectfont\linespread{1}\selectfont\bfseries La répartition de l'hydrogène atomique (\HI) cinématiquement froid et chaud et son lien avec la formation stellaire dans les galaxies proches}}
	\\[1em]
	\normalsize{Présenté le 12 Novembre 2026 par:}\\
	\begin{minipage}{3.5cm}
		\begin{figure}[H]
			\centering
			\includegraphics[width=\linewidth, trim={0 140 0 10cm}, clip]{images/mypic1.jpg}
		\end{figure}
	\end{minipage}\\[10pt]
	\large\textbf{Anjia Randia RABARISOA}\\
	%\large\textbf{Cirno}\\[15pt]
\end{center}
\noindent \flushleft Devant la commission d'examen composée par :

\begin{center}
	\begin{tabular*}{\linewidth}{@{\extracolsep{\fill}}lll}
		Président :    & \textbf{M. Minoson RAKOTOMALALA}        & Professeur Titulaire\\
		Rapporteur :   & \textbf{M. Roger IANJAMASIMANANA}       & Chercheur Postdoctoral\\
		Examinateurs : & \textbf{M. Solohery RANDRIAMAMPANDRY}   & Professeur\\
		& \textbf{M. Toky RANDRIAMAMPANDRY}       & Maître de Conférences\\
	\end{tabular*}
\end{center}
\newpage
\null
\thispagestyle{empty}
\let\cleardoublepage\clearpage